\section{Problem Definition}\label{sec:problem_sup}
We will use $\mc{B}(t, r) := \{\bx \in \R^t\,\mid\,\|\bx\|_2 \leq r\}$  to denote the $\ell_2$-ball in $\R^t$ centered at $\mb{0}$ of radius $r$. 

{\bf Supervised Learning Dataset Distillation.} 
We consider regression tasks over $d$-dimensional real-vectors and real-valued labels. For an $n$-point dataset $D \in (\R^d\times \R)^n$, and a predictor $f : \R^d \to \R$, the mean squared error (MSE) of $f$ on $D$ is $L_{\tn{mse}}(D, f) := (1/n)\sum_{(\bx,y)\in D}(y - f(\bx))^2$. We assume that the training datapoints are norm bounded by $B$ and labels have magnitude at most $b$ for some parameters $B$ and $b$, and impose the same restriction on the synthetic dataset.
Let $D^{\sf sup}_{\tn{train}} = \{(\bx_i, y_i) \in \mc{B}(d,B) \times [-b,b]\}_{i=1}^n$ be an $n$-sized training dataset, while we denote an $m$-sized synthetic dataset by $D^{\sf sup}_{\tn{syn}} = \{(\bz_i, \hat{y}_i)\in \mc{B}(d,B) \times [-b,b]\}_{i=1}^m$. 
By appending a $1$-valued coordinate to feature-vectors we can omit the constant offset and restrict ourselves to linear regressors of the form $\bv^{\sf T}\bx$ as a prediction for the label $y$. We impose $\|\bv\|_2 \leq 1$ as a bound on the norm, and let $\mc{F}_0$ denote the class of such regressors. %The squared loss at $(\bx, y)$ is given by $(y - r^{\sf T}\bx)^2$. \\
We define the \textit{supervised regression dataset distillation} problem as: for a parameter $\eps$, given $D^{\sf sup}_{\tn{train}}$, compute $D^{\sf sup}_{\tn{syn}}$ such that
\begin{equation}
    L^{\tn{err}}_{\tn{mse}}(D^{\sf sup}_{\tn{train}}, D^{\sf sup}_{\tn{syn}}, f) := \left(L_{\tn{mse}}(D^{\sf sup}_{\tn{train}}, f) - L_{\tn{mse}}(D^{\sf sup}_{\tn{syn}}, f)\right)^2 \leq \eps, \qquad  \tn{for all }f \in \mc{F}_0. \label{eqn:errMSEloss}  
\end{equation}


{\bf Offline RL Dataset Distillation.}
Consider a Markov Decision Process (MDP) given by $\langle \mc{S}, \mc{A}, \mc{P}, \mc{R}, \gamma\rangle$, where (i) $\mc{S}$ is the set of states, (ii) $\mc{A}$ is the set of actions, (iii) $\mc{P}$ is the transition probability i.e., $\mc{P}(s'\mid s,a)$ denotes the probability of transitioning to state $s'$ from state $s$ on action $a$, (iv) $\mc{R}$ is the reward function, where $\mc{R}(s,a) \in [0, R_{\tn{max}}]$ is the reward obtained on action $a$ at state $s$, and (v) $\gamma$ is the discount factor. A \emph{policy} $\pi$ is a mapping from states to action, and the goal is to maximize at each state its \emph{value function}: $v_\pi(s) = \E_{\pi}\left[\sum_{t=0}^t \gamma^t R_t\right|s_0 = s]$ where $R_t$ is the reward at step $t$ under the policy starting from state $s$. The action-value function is the expected sum of discounted reward starting from a state and a specific action i.e., $q_{\pi}(s,a) = \E_{\pi}\left[\sum_{t=0}^t \gamma^t R_t\,\mid\, s_0 = s, a_0 = a\right]$. On the other hand, given an action value function predictor $f : \mc{S} \times \mc{A} \to \R$, the corresponding greedy policy $\pi'$ given by $\pi'(s)\in \tn{argmax}_a f(s, a)$ always yields at least as much expected discounted reward starting from any state as $\pi$. 
In offline RL, a dataset $D^{\sf orl}$ consists of a collection  of state, action, reward and next state tuples of the form $(s,a,r,s')$, with generated by some (non-optimal) policies. The goal is to learn a policy from this dataset which maximizes the value function. We cast this problem as deriving a action-value predictor from $D^{\sf orl}$, and taking the greedy policy with respect to it. Under reasonable assumptions on the  MDP and $D^{\sf orl}$, the performance of an action-value predictor $f$ is measured by the Bellman loss (see Appendix \ref{app:misc1} for details) given by:
\begin{equation}
    L_{\tn{Bell}}(D^{\sf orl}, f) = \E_{(s, a, r, s') \leftarrow D^{\sf orl}}\left[\left(f(s, a) - r - \gamma \max_{a' \in \mc{A}} f(s',a')\right)^2\right]. 
\end{equation}
{\it Feature Map and linear action-value predictors.} We assume that $\mc{S},  \mc{A} \subseteq \mc{B}(d_0, B_0)$, and that $\phi : \mc{B}(d_0, B_0)^2 \to \R^d$ is a given feature-map s.t. $\|\phi(s,a)\|_2 \leq B$ for any $(s,a)$ in its domain, for some parameters $B_0$, $d_0$ and $B$. An action-value predictor $f$ is given by $f(s, a) := \bv^{\sf T}\phi(s,a)$, for $(s,a) \in \mc{B}(d_0, B_0)^2$, and we restrict ourselves to the class of such predictors $\mc{Q}_0$ satisfying $\|\bv\|_2 \leq 1$. \\
{\it Datasets.} The training dataset is $D^{\sf orl}_{\tn{train}} = \{(s_i, a_i, r_i, s_i')\}_{i=1}^n \subseteq \mc{S}\times \mc{A} \times [0, R_{\tn{max}}] \times \mc{S}$ of  state, action, reward and next state tuples. Since $\mc{S}$ or $\mc{A}$ could either be discrete or non-convex, for tractable optimization the synthetic dataset $D^{\sf orl}_{\tn{syn}}$ is allowed to consist of tuples $\{(\hat{s}_i, \hat{a}_i, \hat{r}_i, \hat{s}_i') \in \mc{B}(d_0, B_0) \times \mc{B}(d_0, B_0) \times [0, R_{\tn{max}}] \times \mc{B}(d_0, B_0)$. With this, $L_{\tn{Bell}}$ is can defined for $D^{\sf orl}_{\tn{syn}}$ and any $f \in \mc{Q}_0$ as:
\begin{equation}
    L_{\tn{Bell}}(D^{\sf orl}_{\tn{syn}}, f) = \E_{(\hat{s}, \hat{a}, \hat{r}, \hat{s}') \leftarrow D^{\sf orl}_{\tn{syn}}}\left[\left(f(\hat{s}, \hat{a}) - \hat{r} - \gamma \max_{\hat{a}' \in  \mc{A}} f(\hat{s}',\hat{a}')\right)^2\right]. 
\end{equation}
Note that in the above, the maximization inside the loss is taken over the original set of actions, consistent with the definition of the Bellman loss. %which helps optimization in practice for small or finite action spaces. 
With the above setup, we define the \textit{offline RL dataset distillation problem} as follows: For a parameter $\eps$, given $D^{\sf orl}_{\tn{train}}$, compute $D^{\sf orl}_{\tn{syn}}$ such that
\begin{equation}
    L^{\tn{err}}_{\tn{Bell}}(D^{\sf orl}_{\tn{train}}, D^{\sf orl}_{\tn{syn}}, f) := \left(L_{\tn{Bell}}(D^{\sf orl}_{\tn{train}}, f) - L_{\tn{Bell}}(D^{\sf orl}_{\tn{syn}}, f)\right)^2 \leq \eps, \qquad  \tn{for all  }f \in \mc{Q}_0.  \label{eqn:errBell}
\end{equation}


\section{Our Results} \label{sec:our_results}
{\bf Supervised Learning Dataset Distillation.}  For convenience, we shall employ a homogeneous formulation i.e., with $0$ label, using the concatenation $\zeta : \R^d \times \R \to \R^{d+1}$ given by $\zeta(\bx , y) = (x_1, \dots, x_d, y)$ where $\bx = (x_1, \dots, x_d)$.  Observe that $\br^{\sf T}\zeta(\bx, y) = \bv^{\sf T}\bx - y = f(\bx) - y$ where $\br = (v_1, \dots, v_d, -1)$,  and $f \in \mc{F}_0$ s.t. $f(\bx) = \bv^{\sf T}\bx$. Further, if $\|\bv\|_2 \leq 1$ then $1 \leq \|\br\|_2 \leq 2$. Let $\mc{F}$ be the class regressors with target label $0$ where each $h \in \mc{F}$ is given by $h(\bx, y) := \br^{\sf T}\zeta(\bx, y)$ where $1\leq \|\br\|_2 \leq 2$. One can thus extend the notion of the MSE loss $L_{\tn{\tn{mse}}}$ to $\mc{F}$ by letting $L_{\tn{mse}}(D, h) := \E_{(\bx, y)\in D}\left[h(\bx, y)^2\right]$.  We also define $G$ to be a distribution over regressors where a random $g \in G$ is given by sampling $\br \sim N(0,1/(d+1))^{d+1}$ u.a.r. and letting $g(\bx, y) := \br^{\sf T}\zeta(\bx, y)$. With this setup, we prove the following theorem.

\begin{theorem}\label{thm:supervised-distill-1}
    Let $D^{\sf sup}_{\tn{train}}$ the training dataset as described above. For any $\Delta > 0$ and $\delta > 0$, let $g_1, \dots, g_k$  be iid regressors sampled from  $G$ for some $k = O\left({d^2}\log({d}(B+b)/{\Delta})\log(1/\delta)\right)$. Then, with probability $1-\delta$ over the choice of $g_1, \dots, g_k$, if there exists $D^{\sf sup}_{\tn{syn}}$ s.t. 
\begin{equation}
    \sum_{j = 1}^{k} L^{\tn{err}}_{\tn{mse}}(D^{\sf sup}_{\tn{train}}, D^{\sf sup}_{\tn{syn}}, g_j) \leq \Delta', \label{eqn:erronsampledregs}
\end{equation}
then, for all $h \in \mc{F}$, $L^{\tn{err}}_{\tn{mse}}(D^{\sf sup}_{\tn{train}}, D^{\sf sup}_{\tn{syn}}, h) \leq \Delta$, where $\Delta' = \Delta k/O(d^2)$, in particular this holds also for all $f \in \mc{F}_0$. 
\end{theorem}
It can be seen that $L^{\tn{err}}_{\tn{mse}}(D^{\sf sup}_{\tn{train}}, D^{\sf sup}_{\tn{syn}}, g)$ is a convex function over the points of $D^{\sf sup}_{\tn{syn}}$ (see Appendix \ref{app:misc2} for an explanation) and therefore the LHS of \eqref{eqn:erronsampledregs} is also convex and can be minimized efficiently. Based on this we provide the corresponding Algorithm \ref{alg:one}.
\begin{algorithm}

\caption{Supervised Regression Dataset Distillation}\label{alg:one}
{\bf Input:}  $d, k, m \in \mathbb{Z}^+, D^{\sf sup}_{\tn{train}} \in \left(\mc{B}(d, B) \times [-b, b]\right)^n$\\
1. Sample iid at random $g_1, \dots, g_k$ from $G$.\\
2. Output\ \ \ \  $\tn{argmin}_{D^{\sf sup}_{\tn{syn}} \in  \left(\mc{B}(d, B) \times [-b, b]\right)^m} \sum_{j = 1}^{k} L^{\tn{err}}_{\tn{mse}}(D^{\sf sup}_{\tn{train}}, D^{\sf sup}_{\tn{syn}}, g_j)$.
\end{algorithm}

{\bf Lower Bound.} Complementing the above algorithmic result, we prove the following matching (up to logarithmic factors) lower bound on the number of sampled regressors.

\begin{theorem}\label{thm:lbnumber}
For any positive integer $d$, there exists $D^{\sf sup}_{\tn{train}}$ of  $q = (d+1)$ points of the form $\bz = (\bx, y) \in \R^q$ where $\bx \in \R^d, y \in \R$, each of Euclidean norm $\|\bz\|_2 \leq 2$ such that for any choice of homogeneous (i.e., target label $0$) regressors $\{f_t : \R^q \to \R\ \mid\ f_t(\bz) := \bv_t^{\sf T}\bz,  \tn{ where } \bv_t \in \R^q\}_{t=1}^T$,
where $T < q(q+1)/2$, there exists:
\begin{itemize}[noitemsep,nolistsep]
\item $D^{\sf sup}_{\tn{syn}}$ of  $q$ points in $\R^q$ each of Euclidean norm $\leq 2$, and 
\item a regressor $f_0 :\R^q \to \R$ given by $f_0(\bz) := \bv_0^{\sf T}\bz$ for a unit vector $\bv_0 \in \R^q$,
\end{itemize}
satisfying
\begin{equation}\label{eqn:thm:lbnumber}
    L^{\tn{err}}_{\tn{mse}}(D^{\sf sup}_{\tn{train}}, D^{\sf sup}_{\tn{syn}}, f_t) = 0,\ \forall t \in \{1,\dots, T\} \quad \quad \tn{and} \quad \quad L^{\tn{err}}_{\tn{mse}}(D^{\sf sup}_{\tn{train}}, D^{\sf sup}_{\tn{syn}}, f_0) \geq 1/(4q^2).
\end{equation}
 \end{theorem}
 Informally, the above theorem constructs a training dataset such that for any set of less than $q(q+1)/2$ linear regressors one can choose a synthetic dataset on which each of the linear regressors have the same loss as on the training dataset, while there exists a regressor that has significantly different losses. This implies a lower bound of $\Omega(d^2)$ for $k$ in Theorem \ref{thm:supervised-distill-1}.


{\bf Offline RL Dataset Distillation.} For convenience, we define the following modified Bellman loss which incorporates a scale factor $\lambda \in \mathbb{R}$ for the reward in the usual Bellman loss:
\begin{equation}
    \hat{L}_{\tn{Bell}}(D, f, \lambda) := \E_{(s, a, r, s') \leftarrow D}\left[\left(f(s, a) - \lambda r - \gamma \max_{a' \in \mc{A}} f(s',a')\right)^2\right]. \label{eqn:hatLbell}
\end{equation}
To state our result, we need the \emph{pseudo-dimension} ${\sf Pdim}$ (see Appendix \ref{sec:pdim}) of the above loss restricted to single points. In particular, $\mc{U}$ be class of mappings $u: \mc{B}(d_0, B_0)\times \mc{B}(d_0, B_0) \times [0, R_{\tn{max}}]\times \mc{B}(d_0, B_0) \to \R$ where each $u \in \mc{U}$ is defined by $u(s, a, r, s') := \hat{L}_{\tn{Bell}}(\{(s, a, r, s')\}, f, \lambda)$ for some $f \in \mc{Q}_0$ and $\lambda\in [-1,1]$. Let $p := {\sf Pdim}(\mc{U})$. Further, we assume that $\phi$ is $L$-Lipschitz. \\

We also define:
\begin{equation}
    \hat{L}^{\tn{err}}_{\tn{Bell}}(D^{\sf orl}_{\tn{train}}, D^{\sf orl}_{\tn{syn}}, f, \lambda) = \left(\hat{L}_{\tn{Bell}}(D^{\sf orl}_{\tn{train}}, f, \lambda) - \hat{L}_{\tn{Bell}}(D^{\sf orl}_{\tn{syn}}, f, \lambda)\right)^2 \label{eqn:haterrBell}
\end{equation}
We define a distribution $\tilde{H}(\sigma)$ over (predictor, scalar) pairs in which a random $(f, \lambda)$ is sampled by independently choosing $\br \sim N(0,\sigma^2)^d$ and $\lambda \sim N(0,\sigma^2)$ and defining $f(s,a) := \br^{\sf T}\phi(s,a)$. For simplicity we define $H := \tilde{H}(1)$ Using this, we have the following theorem.
\begin{theorem}\label{thm:offlineRL-distill-1}
    Let $D^{\sf orl}_{\tn{train}}$ the offline RL training dataset as described above. For any $\Delta > 0$ and $\delta > 0$, let $(f_1, \lambda_1), \dots, (f_k, \lambda_k)$  be iid samples from  $H$ for some $k = (1/\nu)^{O(d\log d)}O\left(p(\log(1/\nu) + \log(B_0L/(B+R_{\tn{max}})\right)\log(1/\delta)$, where $\nu := \Delta/(B+R_{\tn{max}})^4$. Then, with probability $1-\delta$, if there exists $D^{\sf orl}_{\tn{syn}}$ s.t. 
\begin{equation}
    \sum_{j = 1}^{k} \hat{L}^{\tn{err}}_{\tn{Bell}}(D^{\sf orl}_{\tn{train}}, D^{\sf orl}_{\tn{syn}}, f_j, \lambda_j) \leq \Delta' \label{eqn:erronsampledregsorl}
\end{equation}
then, for all $f \in \mc{Q}_0$, $L^{\tn{err}}_{\tn{Bell}}(D^{\sf orl}_{\tn{train}}, D^{\sf orl}_{\tn{syn}}, f) \leq \Delta$, where $\Delta' = (\Delta/O(1))$. 
\end{theorem}

The following is the distillation algorithm whose guarantees follow directly from Theorem \ref{thm:offlineRL-distill-1}. 
\begin{algorithm}
%\dontprintsemicolon
\caption{Offline RL Dataset Distillation}\label{alg:two}
{\bf Input:}  $d, k, m \in \mathbb{Z}^+$, feature-map $\phi$, $D^{\sf orl}_{\tn{train}} \in \left(\mc{S}\times\mc{A}\times [0, R_{\tn{max}}]\times \mc{S}\right)^n$.\\
1. Sample iid at random $(f_1, \lambda_1) \dots, (f_k, \lambda_k)$ from $H$.\\
2. Output \ \ \ $\tn{argmin}_{D^{\sf orl}_{\tn{syn}} \in  \left(\mc{B}(d_0, B_0) \times \mc{B}(d_0, B_0)\times [0, R_{\tn{max}}] \times \mc{B}(d_0, B_0)\right)^m} \sum_{j = 1}^{k} \hat{L}^{\tn{err}}_{\tn{Bell}}(D^{\sf orl}_{\tn{train}}, D^{\sf orl}_{\tn{syn}}, f_j, \lambda_j)$
\end{algorithm}

{\bf Offline RL with decomposable feature-map.} We consider the natural case when $\phi$ is decomposable i.e., $\phi(s,a) := \phi_1(s) + \phi_2(a) \in \mathbb{R}^d$ for all $s \in \mc{S}, a\in \mc{A}$ for some mappings $\phi_1, \phi_2 : \mc{B}(d_0, B_0)\to \R^d$. Further, we say that $\phi$ is linear and decomposable if $\phi_1$ and $\phi_2$ are linear maps. 


\begin{theorem}\label{thm:offline-rl-decomposable}
 Consider the case when $\phi$ is decomposable as defined above and let $D^{\sf orl}_{\tn{train}}$ the offline RL training dataset. For any $\Delta > 0$ and $\delta > 0$, let $(f_1, \lambda_1), \dots, (f_k, \lambda_k)$  be iid samples from  $\tilde{H}(1/\sqrt{d+1})$ for some $k =O\left({d^2}\log({d}(B+R_{\sf max})/{\Delta})\log(1/\delta)\right)$. Then, with probability $1-\delta$, if there exists $D^{\sf orl}_{\tn{syn}}$ s.t. 
\begin{equation}
    \sum_{j = 1}^{k} \hat{L}^{\tn{err}}_{\tn{Bell}}(D^{\sf orl}_{\tn{train}}, D^{\sf orl}_{\tn{syn}}, f_j, \lambda_j) \leq \Delta' 
\end{equation}
satisfying
\begin{equation}\label{eqn:meancondition}
   \E_{D^{\sf orl}_{\tn{syn}}} [(\phi_1(s), \phi_2(a), r, \phi_1(s'))] = \E_{ D^{\sf orl}_{\tn{train}}} [(\phi_1(s), \phi_2(a), r, \phi_1(s'))] 
   \end{equation}
then, for all $f \in \mc{Q}_0$, $L^{\tn{err}}_{\tn{Bell}}(D^{\sf orl}_{\tn{train}}, D^{\sf orl}_{\tn{syn}}, f) \leq \Delta$, where $\Delta' = \Delta k/O(d^2)$. 
\end{theorem}
In particular, the above theorem shows that the optimization in Algorithm \ref{alg:two} constrained by \eqref{eqn:meancondition}, in the case of decomposable feature-map $\phi$, requires only $\tilde{O}(d^2)$ sampled action-value predictors. Further, when $\phi$ is linear and decomposable i.e., $\phi_1$ and $\phi_2$ are linear maps, then \eqref{eqn:meancondition} is a linear constraint in the points $D^{\sf orl}_{\tn{syn}}$  and it can be shown (see Appendix \ref{sec:orl} for details) that the optimization in Algorithm \ref{alg:two} constrained by \eqref{eqn:meancondition} is convex.



{\bf Discussion of Our Results.} Theorem \ref{thm:supervised-distill-1} and Algorithm \ref{alg:one} together provide an efficient supervised dataset distillation algorithm (for linear regression) with performance guarantees. Specifically, we show that with high probability over $\tilde{O}(d^2)$ sampled regressors, optimizing a convex objective over $D^{\sf sup}_{\tn{syn}}$ to minimize the sum of $L^{\tn{err}}_{\tn{mse}}$ for the sampled regressors, is sufficient to obtain a high quality synthetic dataset. 
While our method adapts the model-training free approach of \cite{ZBDistribMatch} to loss matching, our theoretical guarantees are qualitatively different from those of \cite{nguyen2021dataset} which showed convergence of synthetic features assuming the synthetic labels are given, and of \cite{chen2024provable} who assumed the availability of the optimal trained model. We show that our analysis is tight (upto polylogarithmic factors) by proving a $\Omega(d^2)$ lower bound on the number of sampled regressors in Theorem~\ref{thm:lbnumber}. In addition, we also provide a lower bound of $O(d)$ on the size of the synthetic dataset needed (in Appendix~\ref{sec:lowerbounds}).
We note that Theorem \ref{thm:supervised-distill-1} shows that the randomness (or sample complexity) required is proportional to $\log (1/\Delta)$ where $\Delta$ is the error, which is quantitatively better than sample size proportional to $1/\Delta$ achieved by matrix sketching techniques (see \cite{garg2024distributed}, \cite{wang2018sketched}, \cite{clarkson2009numerical}).
Our theoretical analysis of the supervised regression relies on anti-concentration of $L^{\tn{err}}_{\tn{mse}}$ with respect to random linear regressors, and can be applied to any class of non-linear regressors that admit similar anti-concentration properties.

For offline RL, Theorem \ref{thm:offlineRL-distill-1} and Algorithm \ref{alg:two} give the first dataset distillation algorithm with rigorous performance bounds, without model training. Our novel approach based on matching the Bellman loss w.r.t. sampled $Q$-value predictors differs from the previous behavioral cloning (BC) based methods of \cite{light2024datasetdistillationofflinereinforcement} which matches the loss gradients of the BC objective and of \cite{lei2024offline} which optimizes the synthetic dataset using an action-value weighted BC loss requiring a trained action-value predictor.  

However, the Bellman Loss involves a max term and thus 
 cannot be well represented as a low degree polynomial in the weights of the predictor. Due to this, polynomial anti-concentration (used in the supervised regression case) cannot be applied and the proof is via a conditioning argument. In effect, our algorithm requires sampling $\tn{exp}(O(d\log d)$ predictors in the worst case where $d$ is the output dimension of the feature-map. 

We note that a brute force approach would be to choose a net over the predictors $\mc{Q}_0$ instead of sampling. However, that would necessarily require $\tn{exp}(d)$ such predictors, while in practice a smaller number of randomly sampled predictors suffices, as demonstrated in our experimental evaluations.
\\
While the bound in Theorem \ref{thm:offlineRL-distill-1} is indeed less efficient than $\tilde{O}(d^2)$ sample complexity of the supervised setting, in practice the features are mapped into a smaller dimension than the input mitigating this to some extent. Further, in Theorem \ref{thm:offline-rl-decomposable}, we prove similar $\tilde{O}(d^2)$ in the case of decomposable feature-map $\phi$, and show that the optimization is convex when $\phi$ is linear and decomposable. This is a natural assumption and is the default feature map in applied RL for most value based deep RL approaches. For example, the highly cited work of~\citep{lillicrap2016} explicitly describes concatenating the action and state embeddings. A common trick used in practice and in our experiments is to one-hot encode the actions and concatenate them with the state embedding to obtain the feature map. Additive feature maps are also explicitly studied in~\citep{zhang2020} and~\citep{yang2019}.



{\bf Organization of the Paper.} The proofs of Theorems \ref{thm:supervised-distill-1}, \ref{thm:lbnumber},  \ref{thm:offlineRL-distill-1} and \ref{thm:offline-rl-decomposable} are included in Appendices \ref{sec:Mainpfthm1}, \ref{sec:lbnumber-proof}, \ref{sec:pfthm2} and \ref{sec:orl} respectively. In the next section however, we provide an informal description of the proof techniques. A subset of the experimental evaluations are included in Section \ref{sec:expts} while additional experiments and further details are deferred to Appendix \ref{app:experiments}. 
